\subsection{Adjoint}

PROPOSITION 1. --- Let E and F be two hilbertian spaces. For every mapping \(u \in \mathcal{L}(E; F)\) , there exists a unique mapping \(u^{*} \in \mathcal{L}(F; E)\) such that

\[
\langle u (x) | y \rangle_ {\mathrm{F}} = \langle x | u ^ {*} (y) \rangle_ {\mathrm{E}}\tag{1}
\]

for all \(x \in E\) and all \(y \in F\) . The mapping \(u \mapsto u^{*}\) from \(\mathcal{L}(\mathrm{E}, \mathrm{F})\) into \(\mathcal{L}(\mathrm{F}; \mathrm{E})\) is bijective, isometric and semi-linear (with respect to the automorphism \(\xi \mapsto \bar{\xi}\) of K).

Let \(\mathcal{S}(\mathrm{E},\mathrm{F})\) be the space of all continuous sesquilinear forms on \(\mathbf{E}\times \mathbf{F}\) , endowed with the norm

\[
\| \Phi \| = \sup _ {\| x \| \leqslant 1, \| y \| \leqslant 1} | \Phi (x, y) |.\tag{2}
\]

We define the space \(\mathcal{S}(\mathrm{F},\mathrm{E})\) similarly. We defined (V, p.~16, cor. 2) a Banach space isomorphism from \(\mathcal{L}(\mathrm{E};\mathrm{F})\) onto \(\mathcal{S}(\mathrm{F},\mathrm{E})\) , denoted by \(u\mapsto\Phi_{u}\) and characterized by

\[
\Phi_ {u} (y, x) = \langle y | u (x) \rangle_ {\mathrm{F}} (x \in \mathrm{E}, y \in \mathrm{F}).\tag{3}
\]

In an analogous way we define an isomorphism from \(\mathcal{L}(\mathrm{F},\mathrm{E})\) onto \(\mathcal{S}(\mathrm{E},\mathrm{F})\) . Finally we define a mapping \(\Phi \mapsto \Phi^{*}\) from \(\mathcal{S}(\mathrm{F},\mathrm{E})\) onto \(\mathcal{S}(\mathrm{E},\mathrm{F})\) by

\[
\Phi^ {*} (x, y) = \overline {{\Phi (y , x)}} \quad (x \in \mathrm{E}, y \in \mathrm{F}).\tag{4}
\]

This mapping is bijective, semi-linear and isometric. But formula (1) translates as \(\Phi_{u^{*}} = (\Phi_{u})^{*}\) , hence the proposition.

DEFINITION 1. --- Let E and F be two hilbertian spaces. For every continuous linear mapping \(u: \mathbf{E} \to \mathbf{F}\) , the continuous linear mapping from F into E defined by formula (1) is called the adjoint of \(u\) and is denoted by \(u^*\) .

We have

\[
(u + v) ^ {*} = u ^ {*} + v ^ {*}\tag{5}
\]

\[
(\lambda u) ^ {*} = \overline {{{{\lambda}}}} u ^ {*}\tag{6}
\]

\[
(u ^ {*}) ^ {*} = u\tag{7}
\]

\[
(1 _ {\mathrm{E}}) ^ {*} = 1 _ {\mathrm{E}}\tag{8}
\]

\[
(w u) ^ {*} = u ^ {*} w ^ {*};\tag{9}
\]

in all these formulas, u and v belong to \(\mathcal{L}(\mathrm{E};\mathrm{F})\) , \(\lambda\) is in K, and w in \(\mathcal{L}(\mathrm{F};\mathrm{G})\) where G is a hilbertian space. Formulas (5) and (6) mean that \(u\mapsto u^{*}\) is semi-linear. Formula (8) is obvious. To prove (7), we take the conjugate of the two members of (1), which gives \(\langle u^{*}(y)|x\rangle = \langle y|u(x)\rangle\) , and this proves that u is the adjoint of \(u^{*}\) . Finally, with the notations of (9), we have, for all \(z \in G\)

\[
\langle w (u (x)) | z \rangle = \langle u (x) | w ^ {*} (z) \rangle = \langle x | u ^ {*} (w ^ {*} (z)) \rangle ,
\]

hence \(u^{*}w^{*}\) is the adjoint of wu.

Let \(u: \mathrm{E} \to \mathrm{F}\) be a bijective and continuous linear mapping; then it is also bicontinuous (I, p.~19, cor. 1). From (8) and (9) we immediately deduce that \(u^*\) is bijective and bicontinuous and that

\[
(u ^ {- 1}) ^ {*} = (u ^ {*}) ^ {- 1}.\tag{10}
\]

PROPOSITION 2. --- For every \(u \in \mathcal{L}(\mathrm{E};\mathrm{F})\) , we have

\[
\| u ^ {*} u \| = \| u u ^ {*} \| = \| u \| ^ {2} = \| u ^ {*} \| ^ {2}.\tag{11}
\]

By prop. 1, \(\| u^{*}\| = \| u\|\) , hence \(\| u^{*}u\| \leqslant \| u^{*}\| \cdot \| u\| \leqslant \| u\|^2\) . On the other hand,

\[
\| u \| ^ {2} = \sup _ {\| x \| \leqslant 1} \| u (x) \| ^ {2} = \sup _ {\| x \| \leqslant 1} \langle u (x) | u (x) \rangle = \sup _ {\| x \| \leqslant 1} \langle x | u ^ {*} u (x) \rangle \leqslant \| u ^ {*} u \|,
\]

hence \(\| u^{*}u\| = \| u\|^2\) . Replacing \(u\) by \(u^{*}\) , we get \(\| uu^{*}\| = \| u^{*}\|^2\) , hence (11) follows since \(\| u\| = \| u^{*}\|\) .

Let \(E_{1}, \ldots, E_{n}\) and \(F_{1}, \ldots, F_{n}\) be hilbertian spaces, and for every integer i between 1 and n, let \(u_{i}\) be a continuous linear mapping from \(E_{i}\) into \(F_{i}\) . Then

\[
(u _ {1} \hat {\otimes} _ {2} \dots \hat {\otimes} _ {2} u _ {n}) ^ {*} = u _ {1} ^ {*} \hat {\otimes} _ {2} \dots \hat {\otimes} _ {2} u _ {n} ^ {*}.\tag{12}
\]

Let \(v\) be the continuous linear mapping \(u_1 \hat{\otimes}_2 \ldots \hat{\otimes}_2 u_n\) from

\[
\mathbf {E} = \mathbf {E} _ {1} \hat {\otimes} _ {2} \dots \hat {\otimes} _ {2} \mathbf {E} _ {n} \quad \text {into} \quad \mathbf {F} = \mathbf {F} _ {1} \hat {\otimes} _ {2} \dots \hat {\otimes} _ {2} \mathbf {F} _ {n}
\]

and w the continuous linear mapping \(u_{1}^{*}\hat{\otimes}_{2}\ldots\hat{\otimes}_{2}u_{n}^{*}\) from F into E. It is enough to prove the equality \(\langle y|v(x)\rangle=\langle w(y)|x\rangle\) for \(x\in E\) and \(y\in F\) . By linearity and continuity, we reduce to the case when x and y have the following form

\[
x = x _ {1} \otimes \ldots \otimes x _ {n}, y = y _ {1} \otimes \ldots \otimes y _ {n}
\]

with \(x_{i} \in E_{i}\) and \(y_{i} \in F_{i}\) for \(1 \leqslant i \leqslant n\) . From the definition of scalar product in a tensor product (V, p.~27, formula (6)), we then get

\[
\langle y | v (x) \rangle = \prod_ {i = 1} ^ {n} \left\langle y _ {i} \mid u _ {i} \left(x _ {i}\right) \right\rangle = \prod_ {i = 1} ^ {n} \left\langle u _ {i} ^ {*} \left(y _ {i}\right) \mid x _ {i} \right\rangle = \left\langle w (y) \mid x \right\rangle .
\]

This proves our assertion.

Let E and F be two hilbertian spaces, \(u \in \mathcal{L}(\mathrm{E}; \mathrm{F})\) and n a positive integer. If we put \(u_{1} = \cdots = u_{n} = u\) in formula (12) we obtain the result that the continuous linear mapping \(\hat{\mathbf{T}}^{n}(u^{*})\) from \(\hat{\mathbf{T}}^{n}(\mathrm{F})\) into \(\hat{\mathbf{T}}^{n}(\mathrm{E})\) is the adjoint of the continuous linear mapping \(\hat{\mathbf{T}}^{n}(u)\) from \(\hat{\mathbf{T}}^{n}(\mathrm{E})\) into \(\hat{\mathbf{T}}^{n}(\mathrm{F})\) . The formulas

\[
\hat {\mathbf {S}} ^ {n} (u) ^ {*} = \hat {\mathbf {S}} ^ {n} (u ^ {*}), \quad \hat {\symbf {\Lambda}} ^ {n} (u) ^ {*} = \hat {\symbf {\Lambda}} ^ {n} (u ^ {*})\tag{13}
\]

can be established in the same way as formula (12), on account of the definition of the scalar product in \(\hat{\mathbf{S}}^n (\mathrm{E})\) (V, p.~30, formula (15)) and in \(\hat{\Lambda}^n (\mathrm{E})\) (V, p.~33, formula (26)).

Remark 1. --- Suppose the hilbertian space E does not reduce to 0. We identify \(\mathcal{L}(\mathbf{K};\mathbf{E})\) with E by the mapping \(u\mapsto u(1)\) ; in other words, the vector \(x\) of E is identified with the mapping \(\lambda\mapsto\lambda.x\) from K into E. Then the adjoint of \(x\) is the mapping \(x^{*}:E\to K\) given by \(x^{*}(y)=\langle x|y\rangle\) . In other words, \(x\mapsto x^{*}\) is the canonical semi-linear mapping from E onto its dual (V, p.~15).

Similarly, we identify the number \(\lambda \in \mathbf{K}\) with the endomorphism \(\lambda. 1_{\mathrm{E}}\) of E. Then \(\lambda^{*}\) is precisely the conjugate of \(\lambda\) .

With these identifications, we can define a product \(t_{1} \ldots t_{n}\) where each \(t_{i}\) is, either a number in K, or a vector in E, or a linear form belonging to \(E'\) , or an element of \(\mathcal{L}(E)\) , provided that there are never two consecutive factors \(t_{i}\) and \(t_{i+1}\) of one of the following types :

\(\bullet\) \(xy\) where \(x, y\) are both in \(\mathbf{E}\) , or both in \(\mathbf{E}'\) ;

\(\bullet\) \(x\mathbf{A}\) or \(\mathbf{A}x'\) with \(\mathbf{A} \in \mathcal{L}(\mathbf{E})\) , \(x \in \mathbf{E}\) and \(x' \in \mathbf{E}'\) .

We have the following rules of composition :

\begin{enumerate}
\def\labelenumi{\alph{enumi})}
\item
  associativity;
\item
  every element of \(\mathbf{K}\) commutes with all the other factors;
\item
  we have \((t_1 \ldots t_n)^* = t_n^* \ldots t_1^*\) ; in other words, the adjoint of a product is the product of the adjoints taken in the reverse order. Also \(t^{**} = t\) .
\end{enumerate}

For example, let \(x, y\) be in E and let A be in \(\mathcal{L}(\mathrm{E})\) . Then \(x^{*}y\) represents the scalar product \(\langle x|y\rangle\) and \(x^{*}\mathrm{A}y\) represents the scalar product \(\langle x|\mathrm{A}y\rangle\) . We also have \((\mathrm{A}^{*}x)^{*} = x^{*}\mathrm{A}^{**} = x^{*}\mathrm{A}\) , hence \((\mathrm{A}^{*}x)^{*}y = x^{*}\mathrm{A}y\) , which can be interpreted as

\[
\langle \mathrm{A} ^ {*} x | y \rangle = \langle x | \mathrm{Ay} \rangle
\]

in conformity with the definition of the adjoint. We observe that \(yx^{*}\) is the endomorphism \(z \mapsto y\langle x|z\rangle\) of E, since \(yx^{*}z\) can be interpreted as \(y(x^{*}z)\) by associativity, or as \(y.\langle x|z\rangle\) .

Following Dirac \(^{1}\) , in most works of Mathematical Physics, the elements of E are represented by the symbol \(|x\rangle\) , those of E' by \(\langle t|\) . The scalar product is written as \(\langle x|y\rangle = \langle x|.|y\rangle\) and the first rule of interdiction in the products excludes the combinations of the signs \(\rangle|\) and \(|\langle\) , for example \(|x\rangle|y\rangle\) .

\(^{1}\) See P. A. M. DIRAC, \emph{Quantum Mechanics}, Oxford University Press, New York, 1935.

PROPOSITION 3. --- Let E and F be two hilbertian spaces and \(u \in \mathcal{L}(\mathrm{E};\mathrm{F})\) . The following conditions are equivalent :

\begin{enumerate}
\def\labelenumi{(\roman{enumi})}
\item
  \(u\) is a topological vector space isomorphism, with an inverse equal to \(u^{*}\) ;
\item
  \(u\) is surjective and \(u^{*}u = 1_{\mathrm{E}}\) ;
\item
  \(u\) is injective and \(uu^{*} = 1_{\mathrm{F}}\) ;
\item
  \(u\) is an isomorphism of normed spaces;
\item
  \(u\) is a hilbertian space isomorphism.
\end{enumerate}

Condition (1) means that we have \(u^{*}u = 1_{\mathrm{E}}\) and \(uu^{*} = 1_{\mathrm{F}}\) . Hence the equivalence of (i), (ii) and (iii) follows from E, II, § 3, No.~8, prop. 8. We have already seen the equivalence of (iv) and (v) (V, p.~5). Finally, the relation \(u^{*}u = 1_{\mathrm{E}}\) is equivalent to \(\langle x|u^{*}u(y)\rangle = \langle x|y\rangle\) , that is, to \(\langle u(x)|u(y)\rangle = \langle x|y\rangle\) for all \(x,y\) in E, and evidently implies that \(u\) is injective; this proves the equivalence of (ii) and (v).

An automorphism of the hilbertian space E is also called a unitary operator, that is, an operator \(u \in \mathcal{L}(\mathrm{E})\) satisfying \(uu^{*} = u^{*}u = 1_{E}\) .

Remark 2. --- The relation \(u^{*}u = 1_{\mathrm{E}}\) does not characterize all the automorphisms of the hilbertian space E. For example, let \(\mathrm{E} = \ell^2 (\mathbf{N})\) and let \(u\) be defined by \(u(x_{n}) = x_{n - 1}\) for \(n\geqslant 1\) and \(u(x)_0 = 0\) . We have \(\| u(x)\| = \| x\|\) for all \(x\in \mathrm{E}\) , that is, \(u^{*}u = 1_{\mathrm{E}}\) , but \(u\) is not surjective.

Remark 3. --- The definition (1) of the adjoint \(u^{*}\) can also be written as

\[
\langle y | u (x) \rangle = \langle u ^ {*} (y) | x \rangle ,
\]

or, by V, p.~15, as

\[
\langle u (x), y ^ {*} \rangle = \langle x, (u ^ {*} (y)) ^ {*} \rangle .
\]

But we also have \(\langle u(x), y^* \rangle = \langle x, {}^t u(y^*) \rangle\) , hence we can express the adjoint in terms of the transpose,

\[
(u ^ {*} (y)) ^ {*} = ^ {t} u (y ^ {*}).
\]

